Um die Antikommutatorrelation der Dirac-Matrizen vollständig herzuleiten, beginnen wir mit der Definition des Antikommutators und zeigen, dass er die Clifford-Algebra erfüllt. Die Dirac-Matrizen \(\gamma^\mu\) erfüllen die Relation der Clifford-Algebra:

\begin{align*}
\{\gamma^\mu, \gamma^\nu\} &= \gamma^\mu \gamma^\nu + \gamma^\nu \gamma^\mu = 2 g^{\mu \nu} I.
\end{align*}

Diese Relation folgt direkt aus der Definition der Clifford-Algebra, die besagt, dass das Quadrat einer Dirac-Matrix proportional zur Metrik ist:

\begin{align*}
(\gamma^\mu)^2 &= g^{\mu \mu} I.
\end{align*}

Für den Fall \(\mu \neq \nu\) ergibt sich durch die Anwendung der Antikommutatorrelation:

\begin{align*}
\gamma^\mu \gamma^\nu + \gamma^\nu \gamma^\mu &= 2 g^{\mu \nu} I.
\end{align*}

Damit ist die Antikommutatorrelation vollständig hergeleitet.

Nun zur Hermiteschen Konjugation der Dirac-Matrizen. Die Dirac-Matrizen erfüllen:

\begin{align*}
\gamma^{0 \dagger} &= \gamma^0, \\
\gamma^{i \dagger} &= -\gamma^i \quad \text{für } i = 1, 2, 3.
\end{align*}

Betrachten wir die Hermitesche Konjugation von \(\gamma^\mu\):

\begin{align*}
(\gamma^\mu)^\dagger &= \gamma^0 \gamma^\mu \gamma^0.
\end{align*}

Diese Beziehung folgt aus der Definition der Dirac-Matrizen und ihren Konjugationseigenschaften.

%CHECK: Die Herleitung der Antikommutatorrelation wurde explizit gezeigt. Die Hermitesche Konjugation wurde konsistent hergeleitet.